% 题干转写：完成
% 解析转写：完成
% 图形核对：完成
\section{第35届复赛第1题：地球引力场中的发射运动}
\subsection{题干}
假设地球为质量分布各向同性的球体。空间站位于地面上方高度 $h$ 处，从站内向与地心连线垂直的方向发射小物体。已知地球半径 $R$、质量 $M$、引力常量 $G$；忽略地球自转与大气影响。
\begin{enumerate}
\item 若物体能绕地球作周期运动，初速度大小应满足什么条件？
\item 若初速度为 $v_0$ 且物体会落到地面，求落地速度大小、速度与水平分量的夹角，以及由发射至落地的时间。
\end{enumerate}

\subsection{参考解答整理}
\paragraph{（1）不落地的周期运动条件。} 初始能量和角动量为
\[
E=\frac12mv_0^2-\frac{GMm}{R+h},\qquad L=mv_0(R+h).
\]
有界运动要求 $E<0$。若椭圆轨道恰好与地面相切，切点速度记为 $v$，则
\[
mv_{0\min}(R+h)=mvR,
\qquad
\frac12mv_{0\min}^2-\frac{GMm}{R+h}
=\frac12mv^2-\frac{GMm}{R}.
\]
解得临界初速度
\[
v_{0\min}=\sqrt{\frac{2GMR}{(R+h)(2R+h)}}.
\]
故物体能作不与地面相交的圆或椭圆轨道运动的条件为
\[
\boxed{\sqrt{\frac{2GMR}{(R+h)(2R+h)}}<v_0<\sqrt{\frac{2GM}{R+h}}.}
\]

\paragraph{（2）落地速度、方向与时间。} 当
\[
v_0<\sqrt{\frac{2GMR}{(R+h)(2R+h)}}
\]
时，轨道与地球表面相交。由机械能守恒，落地速度大小为
\[
v=\sqrt{v_0^2+\frac{2GMh}{R(R+h)}}.
\]
由角动量守恒，落地瞬间的水平分量为
\[
v_\theta=\frac{R+h}{R}v_0.
\]
若 $\alpha$ 为速度与水平分量的夹角，则
\[
\alpha=\arccos\!\left(
\frac{(R+h)v_0}{R\sqrt{v_0^2+\frac{2GMh}{R(R+h)}}}
\right).
\]
以极坐标 $(r,\theta)$ 描述运动，守恒关系为
\[
r^2\frac{d\theta}{dt}=(R+h)v_0,
\]
\[
\frac{dr}{dt}=-\sqrt{-\frac{(R+h)^2v_0^2}{r^2}+\frac{2GM}{r}+v_0^2-\frac{2GM}{R+h}}.
\]
因此下落时间由下列定积分唯一确定：
\[
t=\int_R^{R+h}\frac{r\,dr}
{\sqrt{\left(v_0^2-\frac{2GM}{R+h}\right)r^2+2GMr-(R+h)^2v_0^2}}.
\]
令 $a=-(R+h)^2v_0^2$、$b=2GM$、$c=v_0^2-2GM/(R+h)$，将题给积分公式代入端点即可得到官方解答中的闭式时间表达式；该式在临界速度处取相应极限。

\subsection{转写依据}
题干依据 \texttt{CPHO35-Q01.md}，公式依据 \texttt{CPHO35-Q01\_formula\_check\_v1.csv} 与答案页第6至9页。原题无题图，故本题无TikZ图件。
